\chapter{State of the Art}\label{state-of-the-art}

This chapter surveys the work the thesis builds on and, just as importantly, where that work stops. It moves from the classifiers and datasets that define the task, through the large language models and retrieval methods that add context, to the agentic and Big Data ideas that make real-time moderation possible, and closes by naming the gap the thesis sets out to fill. For a recent survey of hate-speech moderation with large models, which maps this same territory in more breadth than a single chapter can, see \citet{hatesurvey2024}.

\section{Detecting hate speech: datasets and a definitional problem}\label{detecting-hate-speech-datasets-and-a-definitional-problem}

The modern study of automated hate-speech detection is shaped by a problem that is as much conceptual as technical. Davidson and colleagues showed that hateful language and merely offensive language draw on the same vocabulary, so any system that treats a word list as a proxy for hate will over-flag ordinary rude speech \citep{davidson2017}. Their three-way split (hate, offensive, neither) became a template for the field precisely because the binary framing hides the cases that matter most.

Later datasets attacked the two hardest parts of the problem. HateXplain added human rationales, marking \emph{which} spans drove an annotator's decision, so that models could be trained and audited for \emph{why} they label something, not only \emph{whether} \citep{hatexplain2021}. ToxiGen went after implicit and adversarial hate, generating large volumes of subtle, machine-written examples that never use an obvious slur \citep{toxigen2022}. Together these three corpora (the ones this thesis fine-tunes on) cover informal slang, explainability, and implicit toxicity.

What they share is also their limit. Each is a fixed, labelled snapshot, and a model trained on them inherits the annotators' context, not the deployment's. When the domain shifts (a gaming stream, a political broadcast, a product review) the same message can flip class, and a statically trained classifier has no way to know. The data got richer; the blindness to live context remained, and that unshaken blindness is the first limitation the design answers: it is why the pipeline does not stop at a statically trained classifier but places a context-aware second tier behind it.

\section{Transformer classifiers: fast, accurate, and context-blind}\label{transformer-classifiers-fast-accurate-and-context-blind}

Transformer encoders are the default tool for text classification, and the family shares one recipe, differing mainly in how hard each pushes it. BERT set the template: bidirectional pretraining that a light head can be fine-tuned on cheaply \citep{bert2019}, with tooling that made such fine-tuning routine \citep{transformers2020}. RoBERTa trades extra pretraining compute for higher accuracy on the same architecture \citep{roberta2019}, while DeBERTa's disentangled attention buys a further step in benchmark scores at the price of more architectural complexity \citep{deberta2021}. HateBERT strikes a different bargain, keeping BERT's design but re-pretraining on abusive English to gain a strong in-domain prior on exactly the task this thesis addresses, at the cost of generality \citep{hatebert2021}. What none of them escapes is the same trade-off between capability and cost, and it is on the cost side that the choice for a live stream is actually made.


Accuracy, though, is not the only axis. Real-time chat needs a classifier that keeps up with the stream, and here DistilBERT is the pragmatic choice: a distilled, six-layer model that keeps most of BERT's quality at a fraction of its size and latency \citep{distilbert2019}. That speed is exactly why this thesis uses it for the first pass.

But every model in this family shares one flaw. Their decision boundary is fixed at training time. They cannot loosen for gaming banter or tighten for a political broadcast, they struggle with sarcasm and coded language, and adapting them to a new domain means retraining. Recent work goes so far as to argue that large language models could replace these traditional classifiers for hate detection outright \citep{rethinkhate2025}, though the cost of doing so per message keeps a fast static tier necessary. A better transformer raises the ceiling; it does not remove the blindness. That is the wall this thesis's second tier is built to get over.

\section{Large language models for content moderation}\label{llms-for-moderation}

Large language models changed what ``reading a message in context'' can mean. Rather than train a fresh classifier, one can prompt an LLM to weigh a message against instructions, and the LLM-as-a-judge line of work showed these models agree with human judgement well enough to serve as evaluators in their own right \citep{llmjudge2023}. Purpose-built safety models followed: Llama Guard frames moderation as an instruction-tuned input-output classifier \citep{llamaguard2023}, riding on the broader capability of the Llama 3 family \citep{llama3_2024}, demonstrating how industry-grade pipelines productionize LLMs for safety. The wider architectural blueprints for deploying these models across enterprise trust-and-safety systems are now surveyed in their own right \citep{platformintegrity2025}, though adapting a model to moderation by fine-tuning is known to carry its own data-engineering pitfalls \citep{adaptmoderation2023}.
There is a trade-off, though: cost. A large model spends orders of magnitude more time per message than a distilled classifier, which makes running one synchronously over a live chat stream impractical. The literature's own framing, an LLM as a careful \emph{judge} rather than a bulk \emph{labeller}, points to the resolution this thesis adopts: use the expensive reasoner selectively, on the cases a fast model cannot settle, and keep it off the hot path.
\section{The fallible judge: how far an LLM evaluator can be trusted}\label{llm-as-a-judge}

The previous section took the reliability of an LLM judge largely for granted; a careful state of the art cannot. The same flexibility that makes these models useful evaluators also makes their judgements fragile in ways a moderation pipeline has to reckon with. A recent survey of the LLM-as-a-judge paradigm catalogues both its appeal as a scalable evaluator and the reliability hazards that come with it \citep{llmjudgesurvey2025} , hazards that weigh all the more in moderation, where the decisive cases are the hard, adversarial ones rather than the average example.

The specific failure modes are by now well documented. Large-scale evaluations find judge models only inconsistently reliable, with high agreement that does not translate into validity and that swings sharply from one benchmark to another \citep{judgereliability2026}. The biases are systematic rather than random: judges favour longer answers, are swayed by the order in which options are presented, and rate outputs from their own model family more highly \citep{llmjudge2023}, and further work shows a verdict can move with the language or surface form of a text rather than its content \citep{judgebias2026}. For moderation in particular, Masud and colleagues show that an LLM asked to moderate can be steered by the persona it is given, returning different verdicts on the same message \citep{hatepersonified2024}. Related work finds LLM moderators carry association biases that make them over-sensitive to particular groups \citep{moderationbias2025}, and the judge is itself an attack surface, as a recent security systematisation makes plain \citep{judgesecurity2026}.

None of this makes an LLM judge unusable; it makes an unexamined one dangerous. The strength of the paradigm, flexible and explainable judgement without a fresh training set, is exactly why this thesis builds its second tier on it. The limitation the literature exposes is that reliability is measured almost entirely in offline, single-shot settings on general benchmarks, and rarely for a judge applied selectively to the uncertain slice of a live moderation stream. This thesis narrows that gap in ways that follow directly from the biases above: it pins the judge to a deterministic, temperature-zero setting to blunt surface-form variance, and it treats the judge as one configuration among several, measured head to head against a static baseline and against a multi-agent decomposition, so its verdicts are reported as evidence rather than trusted on faith.

\section{Beyond the verdict: rationale generation and its faithfulness}\label{explainability-rationale}

A verdict a moderator cannot question is a verdict a moderator cannot trust. This matters most for a system whose second tier overturns the first: when the slow judge contradicts the fast classifier, a human reviewer needs to see why before accepting the reversal. The field recognised this early. HateXplain paired every label with human rationales, the specific spans that drove the decision, so that a model could be trained and audited for why it flags something and not only whether \citep{hatexplain2021}. The harder question, which rationales can actually be believed, was formalised by DeYoung and colleagues, whose ERASER benchmark measures an explanation's faithfulness, whether removing the highlighted evidence actually changes the prediction, rather than merely its plausibility to a reader \citep{deyoung2020eraser}.

Large language models shifted explanation from highlighting spans to writing out reasoning. Recent moderation work asks the model to show its steps: Yang and colleagues generate an explicit chain of reasoning alongside the label, so the explanation is the argument rather than a gloss added afterwards \citep{hare2023}, while others use the model to extract human-readable rationales that make an otherwise opaque classifier interpretable \citep{llmrationales2024}. Purpose-built explainable detectors now ship the label together with a generated rationale, whether extracted from the text \citep{targe2025} or grounded in moral-foundations reasoning \citep{smra2026}, though recent work cautions that even the human rationales used as ground truth disagree, which complicates how such explanations ought to be judged \citep{disagreerationales2026}. This suits an LLM judge, which already emits text: the same call that decides can also justify.

The strength of this line for a two-tier design is plain, since an explanation turns an overturned verdict from an unexplained veto into an auditable decision, exactly what a human in the loop needs. The limitation is equally plain, and recent work is blunt about it. Masud and colleagues show that when LLMs act as moderators their explanations can be fluent yet unfaithful, confidently rationalising decisions that shift with the persona or framing they are handed \citep{hatepersonified2024}, and the faithfulness measures behind ERASER are seldom applied to moderation rationales at all, still less to ones produced live over a stream. The gap this thesis occupies is therefore practical rather than metrological. Its Tier-2 judge must return a written justification with every verdict, not as decoration but as the artefact a reviewer reads when deciding whether to trust an overturn, and the multi-agent variant records the intermediate reasoning of each specialist alongside the final decision, so a wrong verdict can be traced to the step that produced it. What this thesis does not do is score those explanations for faithfulness in the ERASER sense; the explanations are used qualitatively, to diagnose how the judge fails (Section~\ref{how-the-judge-fails}), and measuring their faithfulness is left open.

\section{Grounding the judge: retrieval-augmented generation and its newer forms}\label{retrieval-augmented-generation}

One way to make an LLM's judgement less arbitrary is to ground it in retrieved evidence. Retrieval-augmented generation fetches relevant passages from an external store and conditions the model on them, improving factuality and letting the knowledge change without retraining the model \citep{rag2020}. The retrieval half rests on dense sentence embeddings, following Sentence-BERT \citep{sbert2019}, with model choice guided by benchmarks such as MTEB \citep{mteb2023}. On the storage side, retrieval at scale rests on an approximate-nearest-neighbour index over those embeddings, the core technique behind modern vector databases, from the Hierarchical Navigable Small World graphs that back stores such as Qdrant \citep{hnsw2020} to GPU-scale libraries like FAISS \citep{faiss2021}, whose designs and trade-offs are surveyed by Pan and colleagues \citep{vecdbsurvey2024}.

Retrieval has not stood still since. A systematic review of the area traces the progression from naive passage retrieval to advanced, modular pipelines that rerank, fuse, and iterate \citep{ragreview2025}, and the current frontier is agentic retrieval, in which the model itself decides what to fetch, when, and how to use it, rather than following a fixed retrieve-then-read recipe \citep{agenticragsurvey2025}. A parallel line replaces flat passages with structured knowledge: graph-based retrieval, and for moderation specifically a meta-toxic knowledge graph that hands the model related toxic concepts rather than loose text \citep{metatoxicKG2024}. However, building and updating a knowledge graph in real-time over a high-volume chat stream introduces substantial latency and state-management overhead. The limitation, for its purposes, is that this machinery is developed and benchmarked for open-domain question answering, where the goal is factual recall, not for a low-latency moderation decision. The thesis therefore uses vector retrieval in the deliberately simple, evaluable form set out below, and leaves graph-structured retrieval as future work.

Applied to moderation, RAG is naturally read as \emph{precedent}: retrieve past cases that resemble the current message and let them inform the verdict, the way a moderator recalls similar rulings. It is a strong idea, but it inherits the quality of its corpus, and it retrieves \emph{examples} rather than \emph{rules}. A store of past cases can tell the judge ``we saw something like this before''; it is less suited to telling it ``here is the policy that governs this.'' That distinction motivates one of the thesis's own experiments, and it is taken up next.

\section{Knowledge-grounded moderation and the LLM-Wiki}\label{knowledge-grounded-moderation-and-the-llm-wiki}

Not all retrieval targets are alike. Where RAG typically indexes a corpus of raw documents or past examples, a parallel idea is to consult a \emph{curated, structured} knowledge base: policies, definitions, glossaries, and edge-case rulings organised as interlinked pages rather than as loose text. Karpathy's ``LLM-Wiki'' pattern is a recent articulation of this \citep{karpathy2026llmwiki}, in which an LLM reads sources once and maintains a persistent, cross-referenced wiki that compounds over time, rather than re-querying raw documents on every question. Notably, at the scale that pattern is aimed at it needs no retrieval machinery: an index page tells the model which pages to open and they are read straight into context, with search suggested only once a wiki outgrows that.

For moderation the contrast is sharp and, as far as the surveyed literature goes, under-explored: vector RAG retrieves similar past \emph{cases} (bottom-up, example-driven), while a wiki-style knowledge base supplies curated \emph{knowledge} (top-down, rule- and definition-driven), including the coded language and contextual exceptions that rarely surface as clean nearest neighbours. The pieces exist separately: Lewis et al. \citep{rag2020} retrieve document passages by similarity, while safety systems such as Llama Guard \citep{llamaguard2023} classify with no external retrieval at all. Neither weighs a curated wiki-style knowledge backend against vector retrieval for a moderation decision. This thesis treats the two as switchable backends behind the same judge and runs exactly that comparison. Graph-structured retrieval (for example GraphRAG) is a further step along this axis and is noted as future work.

\section{From one prompt to many agents: decomposition and orchestration}\label{agentic-and-multi-agent-systems-for-trust-and-safety}

Beyond a single prompted call, agentic methods let a model reason and act in steps. ReAct interleaves reasoning traces with tool use, so a model can gather evidence before committing to an answer \citep{react2023}. The natural extension for moderation is to decompose the decision across specialists (a risk assessor, a behavioural profiler, an escalation router, a supervisor) rather than ask one prompt to do everything at once.

The single-agent, single-call picture has since grown into a field of its own. Recent surveys of LLM multi-agent orchestration organise the design space by how agents are coordinated, distinguishing centralised, decentralised, and hierarchical topologies, the last of which places a supervisor over specialised workers \citep{maorchestration2026}, while companion work catalogues the architectures and protocols by which such systems hand off work and reconcile conflicting outputs \citep{orchestrationarch2026}. The moderation decomposition this thesis adopts, a risk scorer and a behavioural profiler feeding an escalation router and a final supervisor, is a hierarchical, supervisor-led instance of exactly this pattern. What the orchestration literature adds is both vocabulary and caution: these same surveys report that a large share of multi-agent systems fail in production through mis-specification and unreconciled disagreement, which is precisely why this thesis keeps its chain short, fixed, and measured rather than open-ended. Moderation-specific instances of the pattern are beginning to appear: Gajewska and colleagues, for instance, pair a central moderator agent with dynamically constructed community agents that each contribute a demographic group's perspective, improving implicit-hate detection on ToxiGen \citep{communitymod2026}. Such systems are still evaluated offline on fixed benchmarks and run as standalone detectors, however, rather than as a selective tier behind a fast classifier on a live stream , the setting this thesis targets.

Two limitations follow the specific work. ReAct \citep{react2023} is demonstrated on offline question-answering and interactive-decision benchmarks, not on a live moderation stream, and it keeps the reasoning in a single agent rather than splitting it across specialists. Where a specialist decomposition is proposed for Trust and Safety, it is usually described rather than measured against the monolithic judge it replaces. This thesis addresses both: it runs the agent chain asynchronously on flagged cases so the stream is never blocked, and it evaluates the multi-agent verdict head to head against the single-judge baseline.

\section{The moderator stays in the loop: people and machines together}\label{human-in-the-loop}

Automated moderation does not run in a vacuum; it sits inside a human system, and the literature that takes that system seriously is a necessary counterweight to any claim of full automation. Gillespie frames platform moderation as a set of hidden, consequential human decisions rather than a solved engineering problem \citep{gillespie2018custodians}, and Roberts documents the human labour that industrial-scale moderation actually rests on, the reviewers who see what the classifiers miss \citep{roberts2019behindscreen}. The lesson both draw is that the real design question is not whether to replace human moderators but where to place machine and human so that each does what it is best at.

A concrete line of systems work builds exactly that division of labour. Jhaver and colleagues study Reddit's Automoderator and show how a rule-based bot and human moderators share the load, the machine taking the clear cases and the people the judgement calls \citep{jhaver2019automod}, and Chandrasekharan and colleagues go further with Crossmod, a machine-learning system that surfaces borderline content for human moderators instead of deciding for them \citep{chandrasekharan2019crossmod}. The pattern is consistent: automation as a filter and a triage layer, with the human kept in the loop for the decisions that carry the most risk.

On live-streaming platforms specifically, the domain this thesis targets, that division of labour is carried almost entirely by unpaid volunteer moderators. Wohn's study of Twitch moderators documents who they are and the emotional labour the role demands \citep{wohn2019volunteer}; Seering and Kairam quantify what they actually do and how heavily they lean on the platform's own tools, chiefly Twitch's AutoMod word filter and the ban, timeout, and delete commands, supplemented by third-party bots built on the open moderation API \citep{seering2023who}; and Cai and colleagues show that on larger channels moderation is not a solo task at all but a coordinated team effort, with several moderators dividing real-time chat duty, agreeing rulebooks off-stream, and handing off to one another during a raid or a pile-on \citep{cai2022coordination}. The tooling that supports them today is overwhelmingly rule-based, a blocklist and a set of manual actions, which is precisely the ceiling that an automated, context-aware tier is meant to raise.

The arrival of large language models has sharpened this pattern rather than dissolved it. Park and colleagues build an explicitly collaborative system in which an LLM and human moderators divide cross-cultural hate-speech decisions between them \citep{llmc3mod2025}; Bachar and colleagues learn when an automated tier should hand a case up to a human at all, treating escalation itself as a prediction problem \citep{escalate2026}; and, against the assumption that bigger is always better, Zhan and colleagues find that small specialised models can outperform large ones at moderation \citep{slmmod2024}, a finding this thesis takes literally rather than as a challenge to it: the fine-tuned DistilBERT does the high-volume first pass cheaply, and the large model is reserved only for the selective, context-heavy judgements a small model cannot make.

What this literature settles is a principle worth stating plainly: \textbf{effective moderation is a human--AI collaboration, in which the machine's job is to filter, triage, and escalate, not to have the last word on the hardest cases}. What it mostly stops short of is wiring that principle into a running real-time pipeline with a measured escalation boundary. This thesis takes the principle literally. Its multi-agent tier ends in an escalation router whose available actions include sending a case to human review, which is the Crossmod and Automoderator idea expressed as one component of a streaming system, and its two-tier structure, a cheap fast filter with an expensive reasoner reserved for the uncertain minority, is the cascade this line of work recommends, made concrete and then measured. The human is not simulated in the experiments, but the seam where a human would take over is built, named, and evaluated as part of the pipeline.

\section{No message is an island: conversation and author history}\label{contextual-behavioural}

A single message is a thin basis for a decision. A borderline line from a first-time poster and the same words from a serial offender deserve different treatment, and a reply that reads as an attack in isolation may be ordinary banter given the thread that came before it. This is more than intuition. Pavlopoulos and colleagues found that annotators who are shown the preceding conversation revise their toxicity judgements often enough to matter, and that a classifier given no context systematically misreads exactly the borderline cases where the decision is hardest \citep{pavlopoulos2020context}. The setting itself was charted earlier by Wulczyn, Thain and Dixon, who studied personal attacks at the scale of entire conversations rather than as isolated strings \citep{wulczyn2017exmachina}.

Two strands develop the idea. The first is conversational: what came before shapes what a message means. Zhang and colleagues forecast when a still-civil exchange is about to turn hostile from the conversation's early history alone, evidence that the thread carries a signal any single line hides \citep{zhang2018awry}, and more recent work extends the same logic to implicit hate that only resolves once the surrounding turns are read \citep{ghosh2023cosyn}. The second strand is behavioural: who is speaking matters as much as what is said. Mishra and colleagues build a profile of an author from their earlier posts and show that this history improves abuse detection over content taken in isolation \citep{mishra2018author}.

The most recent work sharpens the framing rather than settling it. Berezin and colleagues argue that toxicity is better measured as contextual harm than as a fixed property of the text, which is precisely the reorientation a context-aware moderator needs \citep{contextualharm2025}, and current large-language-model detectors fold the surrounding context directly into the prompt to catch implicit cases that a context-blind classifier lets through \citep{implicitctx2025}.

What this body of work establishes is, for the design of a second tier, foundational: context and history change the correct answer, and sometimes by a measurable margin. What most of it leaves open is acting on that under live conditions. These studies run on fixed, offline datasets, usually one platform at a time, and they treat context as a static feature bundled in at training time rather than as a signal assembled on demand while a stream flows past. The behavioural half is thinner still, because public corpora rarely follow the same author across enough messages to build a dependable profile. The recommendation to use history is common; a system that assembles a user's recent record, the current thread, and a short behavioural fingerprint at the moment of judgement, injects that bundle as one switchable module, and then measures whether it actually changed the verdict, is not. That gap is what the temporal-memory module of this thesis is built to fill, and its effect, positive or not, is reported in the evaluation chapter.

\section{Moving the stream at scale: Big Data infrastructure for real-time moderation}\label{big-data-streaming-for-real-time-nlp}

The reasoning above is moot if the data cannot be moved. Real-time moderation is a streaming problem, and the standard tools are a durable log for ingestion and a distributed engine for processing: Apache Kafka for high-throughput, replayable queues \citep{kafka2011} and Apache Spark for distributed computation over them \citep{spark2016}, typically arranged as a stream-first (Kappa) pipeline. Kafka's fitness for exactly this high-throughput ingestion role is not taken on trust here: it has been evaluated directly for big-data streaming workloads \citep{kafkaeval2017}, and the trade-offs among competing stream-processing frameworks are themselves a subject of systematic study rather than a matter of taste \citep{streamframeworks2019}. The deeper appeal is scalability: because a partitioned log and a distributed compute engine both grow by adding nodes rather than by rewriting the design, the same architecture that serves one stream serves many, and the storage-centric layout lets new context modules and even new platforms be attached without disturbing the flow beneath them, so the system's capacity expands with the hardware rather than being capped by its structure. Close to this thesis's setting, Poornima and colleagues built a real-time sentiment system over YouTube live chat on exactly this Spark-and-Kafka backbone \citep{poornima2025}, demonstrating the pattern's fit for live social streams. What that line of work establishes is the plumbing; what it does not attempt is to layer selective LLM reasoning and multiple forms of context on top, which is where this thesis extends it.

\section{Evaluating under heavy class imbalance}\label{evaluating-under-heavy-class-imbalance}

A final thread concerns how results are judged. Real moderation traffic is overwhelmingly normal, so a model that never flags anything can score very high accuracy while catching nothing, which makes raw accuracy actively misleading. The honest alternative is to lead with precision, recall, and F1 on the toxic class, and to test whether a difference between two systems is real rather than noise. McNemar's test is the standard paired test for exactly this comparison of two classifiers on the same data \citep{mcnemar1947}, as argued in the classic treatment of statistical tests for learning algorithms \citep{dietterich1998}, with bootstrap confidence intervals a common complement \citep{efron1993}. Evaluation is itself being rethought for LLM-based moderators, with frameworks that fold in socio-cultural context rather than a single global label \citep{socioculturaleval2024}. This thesis adopts that significance-tested, imbalance-aware methodology throughout.

\section{Positioning: the gap this thesis fills}\label{positioning-the-gap-this-thesis-fills}

Read together, these works supply strong parts that no single system assembles, and naming the specific limitation of each is what defines the gap. Davidson and colleagues \citep{davidson2017}, and the transformer classifiers that followed such as HateBERT \citep{hatebert2021}, deliver a fast first pass, but each fixes its decision boundary at training time and so cannot adjust to a live domain without retraining. The LLM-as-a-judge line \citep{llmjudge2023} shows a large model can judge about as reliably as a human, yet more recent evaluation is blunt that this reliability is inconsistent and biased \citep{judgereliability2026, judgebias2026} and is measured offline rather than per message over a stream; Llama Guard \citep{llamaguard2023} brings that judgement to moderation but runs its model on every input in a single pass, with no fast tier to make stream-scale use affordable. The explainability literature can produce a rationale for a verdict \citep{deyoung2020eraser} but warns that LLM explanations are often fluent rather than faithful \citep{hatepersonified2024}. Lewis and colleagues \citep{rag2020} ground a model in retrieved passages, and moderation systems that borrow the pattern retrieve past \emph{cases} as precedent, but none weigh that retrieval against a curated, wiki-style \emph{knowledge} base. ReAct \citep{react2023} interleaves reasoning and action on offline benchmarks, not a live stream, and does not split the safety decision into specialised agents or measure such a split. The human-in-the-loop tradition recommends exactly the machine-triage-with-human-escalation cascade this problem needs \citep{chandrasekharan2019crossmod, escalate2026}, but rarely wires it into a measured real-time pipeline. And Poornima and colleagues \citep{poornima2025} carry real-time analysis on a Kafka-and-Spark backbone, but with a lightweight sentiment model rather than a tiered LLM reasoner.

Each of these limitations maps directly onto a choice in the proposed architecture, and making that mapping explicit is the point of this section. Because a fast classifier is context-blind while a large model is too slow to run on every message, the pipeline splits into a static Tier-1 and a selective Tier-2 behind a confidence gate, so speed and context are not traded against each other. Because context and history change the right answer yet are usually left as an offline recommendation, that context is assembled at judgement time into a temporal-memory module. Because retrieved cases are not the same as curated rules, RAG and an LLM-Wiki are built as switchable backends and compared head to head. Because a single agent leaves the safety decision undifferentiated, the judge is decomposed into a risk scorer, a behavioural profiler, an escalation router and a supervisor, a hierarchical topology whose documented failure modes are the reason the chain is kept short and fixed. Because an LLM judge is unreliable and biased, it is pinned to a deterministic setting and reported as one measured configuration among several rather than trusted on faith. Because explanations can be unfaithful, the judge must return a written justification with every verdict, and the multi-agent chain stores each specialist's reasoning separately, so a verdict can be traced step by step rather than taken on trust. Because effective moderation is a collaboration, the escalation router's available actions include handing a case to a human. And because throughput cannot be an afterthought, the whole pipeline sits on a Kafka-and-Spark streaming backbone whose fitness for the role is established rather than assumed \citep{kafkaeval2017}.

The gap, then, is not any one of these capabilities but their combination, and the honest evaluation of it. No surveyed system decouples a fast static tier from a selective LLM tier, injects temporal, retrieved, and multi-agent context as independently switchable modules, sets vector RAG against a curated-knowledge backend, requires an explanation with every reversal, reserves a route to human review, and measures each configuration under the class imbalance that real traffic imposes. That is what this thesis builds and tests: a single real-time pipeline that (1) keeps a fast static tier and a slow LLM tier apart so context and throughput coexist, (2) injects temporal memory, retrieved precedent, and a multi-agent decomposition as separate, switchable modules, (3) adds a curated-knowledge LLM-Wiki backend and compares it head to head with vector RAG, (4) treats the LLM judge as a fallible component, pinned to determinism, made to explain itself, and left open to human escalation rather than trusted outright, and (5) reports every configuration with imbalance-aware, significance-tested metrics. The chapters that follow describe that system and put each claim to the test.